% Nouveau EGA V (cote 25): archive cote 25, site page 16 (PDF page 17).
% typescript. Draft compared: none.
% AI-transcribed by Claude Opus 5.5 (claude-opus-5-5, adaptive thinking, effort high), Message Batches API, batch ega25b (msgbatch_01A7n5YuXubUPNK1zsepi8RG), 2026-11-13; not yet checked by eye.
\scanpage{25/16}
\mspage{13}

\deleted{\uline{Corollaire} 2.8. Sous les conditions de 2.7. , pour que $Z_i$ soit non
immergé, il f et s que $Z_i$ \inserted[typed]{sur \unsure{$\bar{k}$}} soit non immergé. S'il en est ainsi, la
multiplicité radicielle de $F$ en $Z_i$ est égale à la multiplicité radici-%
elle de $G$ sur $k(\ \ )$ en $Z_i$ .

La première assertion résulte aussitôt de la dernière assertion
de 2.7. Quant à la première, on peut supposer\xout{, (comme il s'agit d'une}
\inserted[typed]{par restriction à un ouvert de X et}
\xout{\unsure{notion géométrique}} par dévissage sur $F$ (comme $Y$ est plat sur $X$
donc $F\ \ G$ est une foncteur exact) que l'on a $F=\underline{O}_{Z_i}$ , et on peut
évidemment remplacer $X$ par $Z_i$ . On est donc ramené au cas où $X$ est intè-%
gre et $F=\underline{O}_X$.} Soit $K/k$ le\xout{s} corps des fonctions de $X$, et soit $n=\dim P$;
introduisant des coordonnées \xout{projec} affines $T_1,..,T_n$ dans $P$ (en choisiss
ant un hyperplan à l'infini $H_\infty$ tel que $f(X)$ ne soit pas contenu dans
celui-ci) et $S_1,...S_n$ dans $\check{P}$ , on voit que le corps\deleted{\unsure{ex}}\inserted{des} fonctions $L$
de $Y_{\hand{\unsure{y}}}$ \xout{xxx} s'identifie au corps des fractions de l'anneau intègre
$K[S_1,..,S_n]/(\sum t_iS_i - 1)$ , où les $t_i \in K$ sont les images des $T_i$ à
l'adde de $f:X\to P$ . Comme $\dim \overline{f(X)} > 0$ , les $t_i$ ne sont pas tous
algébriques sur $k$, à fortiori ils ne sont pas tous nuls, supposons
p. ex. $t_n \neq 0$. On constate alors aussitôt qu'on a $L = K(S_1,..,S_{n-1})$ ,
(extension transcendante pure), $S_n\in L$ étant donné par l'équation
\[ \sum t_iS_i - 1 = 0 \]
en fonction des $S_i$ $(1\leq i\leq n-1)$ et des $t_i$ $(1\leq i\leq n)$. D'autre part,
$k' = k(\hand{y})$ s'identifie à $k(S_1,..,S_n)$, et l\xout{a}'inclusion canonique
$k'\longrightarrow L$ s'obtient en envoyant $S_i$ sur $S_i$ , i.e. $k'$ en tant que sous-%
\xout{\unsure{exten}}\inserted[typed]{extension} de $L$ est la sous-extension engendrée par les $S_i$ $(1\leq i\leq n)$, ou
ce qui revient évidemment au même, par les $S_i$ $(1\leq i\leq \xout{n}\inserted[typed]{n-1}\hand{)}$ et par
$S_n = a_o + a_1S_1 + \ldots + a_{n-1}S_{n-1}$ , où $a_o = t_n^{-1}$ , $a_i =-t_it_n^{-1}$ pour
$1\leq i\leq n-1$. \deleted{Comme}\inserted{Notons que} le corps engendré par les $a_i$ ou par les $t_i$ est évidem
ment le même, \deleted{les $a_i$ \unsure{ne} sont également pas tous algébriques sur k.}\inserted{leur degré de transcendance commun
n'est autre que $\dim f(X)$.}

\margin{16}
